\emph{(i)} Left and right multiplications commute, so the hypercharge commutes with the weak generators; with $T_{a} = \tfrac12B_{a}$, $(B_{1}, B_{2}, B_{3}) = (e_{12}, e_{20}, e_{01})$, $[T_{a}, T_{b}]$ is $\mp T_{c}$ for $(a, b, c)$ cyclic (the bivector algebra of $\Cl(3,0)$, the $\su(2)$ of Theorem \claimref{standard-model.gauge-algebra-minimal-massive-composite} in the bivector realisation of \claimref{synthesis.field-level-brst-differential-of-the-derived-gauge-fields-is-nilpotent}), so $\rho$ is a representation; $\widetilde{B_{a}} = -B_{a}$ and $\widetilde{e_{12}} = -e_{12}$ make every $\rho(X)$ antisymmetric for $\langle\widetilde{\phi}\psi\rangle_{0}$. The check confirms the brackets and the antisymmetry on samples, and that the $32 \times 8$ matrix stacking the four generators has rank $8$, so no nonzero vector is fixed by all of them. With $sc = -\tfrac{g}{2}[c, c]$ and $s\Phi = -g\rho(c)\Phi$, $s^{2}\Phi = -g\rho(sc)\Phi + g\rho(c)s\Phi = \tfrac{g^{2}}{2}\rho([c, c])\Phi - g^{2}\rho(c)\rho(c)\Phi = 0$, because $\rho(c)\rho(c) = \tfrac12\rho([c, c])$ for the odd ghost and $\rho$ is a representation; this holds for every value of $\Phi$, so the full differential needs nothing of the vacuum (the same lemma, on the whole jet algebra, is the intended citation named in the extraction note). The free part: $s_{0}^{2} = 0$ since $s_{0}c = 0$. Split $sS = 0$ by field degree; at degree two, $s_{0}S_{2} + s_{1}S_{1} = 0$. The kinetic term $\tfrac12\langle D_{\mu}\Phi, D^{\mu}\Phi\rangle$ about a constant $v$ has no linear part, and the gauge-fixing and ghost terms are $s\Psi$ with $\Psi = \bar c^{a}(F_{a} + \tfrac{\xi}{2}b_{a})$, $F_{a} = \nabla\cdot A_{a} - \xi g\langle\rho(T_{a})v, h\rangle$, which has no part of degree below two, so at degree two they are $s_{0}\Psi_{2}$, invariant under $s_{0}$. So $S_{1} = \int\langle\nabla V(v), h\rangle$ and $s_{0}S_{2} = -s_{1}S_{1} = g\int\langle\nabla V(v), \rho(c)h\rangle = -g\int c^{a}\langle\rho(T_{a})\nabla V(v), h\rangle$. This vanishes for all $h$ and $c$ exactly when $\rho(X)\nabla V(v) = 0$ for every $X$, and since no nonzero vector is invariant, exactly when $\nabla V(v) = 0$. The same identity in finite form is the derivative of gauge invariance, $\langle\nabla V(\Phi), \rho(X)\Phi\rangle = 0$, along $h$: $\mathrm{Hess}\,V(v)(h, \rho(X)v) + \langle\nabla V(v), \rho(X)h\rangle = 0$. The check computes gradients and Hessians of a quartic exactly (interpolating $t \mapsto V(v + th)$, a polynomial of degree four), confirms this identity at a stationary and a non-stationary point, finds a nonzero obstruction $\langle\nabla V(v), \rho(X)h\rangle$ at the non-stationary one, and a zero gradient at the vacua of (iii).

\emph{(ii)} Rotors act as $\phi_{k} \mapsto U\phi_{k}$ and $\phi_{k} \mapsto \phi_{k}R$ with $R \in \mathrm{span}\{1, e_{12}\}$. Then $\widetilde{U\phi_{1}}U\phi_{2} = \psi$ and $\widetilde{\phi_{1}R}\phi_{2}R = \widetilde{R}\psi R$, which fixes the span$\{1, e_{12}\}$ part of $\psi$ and rotates its span$\{e_{02}, e_{01}\}$ part (these anticommute with $e_{12}$), so $N_{k}$, $K$ and $K'$ are invariant (checked under rational rotors, with a non-gauge map that changes $K$). $K + K' = \langle\widetilde{\psi}\psi\rangle_{0} = N_{1}N_{2}$, the norm of $\Cl(3,0)^{+}$ being multiplicative. The rank of the Goldstone map is $4$ minus the dimension of the stabiliser. Conjugating by a weak rotor, take $\phi_{1}$ scalar when $\phi_{1} \ne 0$. For $X = \sum_{a}x_{a}T_{a} + x_{Y}Y$, $\rho(X)\phi_{1} = \phi_{1}(\tfrac12(x_{1}e_{12} + x_{2}e_{20} + x_{3}e_{01}) + yx_{Y}e_{12})$, which vanishes exactly when $x_{2} = x_{3} = 0$ and $x_{1} = -2yx_{Y}$: the stabiliser of a scalar is spanned by $Z = T_{1} - \tfrac{1}{2y}Y$, acting as $Z\phi = \tfrac12(e_{12}\phi - \phi e_{12})$, and $T_{1}\phi_{1}, T_{2}\phi_{1}, T_{3}\phi_{1}$ are independent. $Z$ kills span$\{1, e_{12}\}$ and acts on span$\{e_{02}, e_{01}\}$ as left multiplication by $e_{12}$, invertibly. If $K' = 0$, then $\psi \in \mathrm{span}\{1, e_{12}\}$, $\phi_{2} = \phi_{1}z$ with $z = \psi/N_{1}$, so $Z\phi_{2} = 0$ and the rank is $3$ (and symmetrically if $\phi_{1} = 0 \ne \phi_{2}$). If $K' \ne 0$, both doublets are nonzero and $\phi_{2} = \phi_{1}\psi/N_{1}$ has a nonzero span$\{e_{02}, e_{01}\}$ part, so $Z\phi_{2} \ne 0$: the stabiliser is trivial and the rank is $4$. The check confirms the rank and the identity on all $3^{8}$ pairs with entries in $\{-1, 0, 1\}$: rank $0$ once, $3$ on $672$ pairs, $4$ on $5888$, each exactly as $K'$ predicts.

\emph{(iii)} With $S = N_{1} + N_{2}$, $D = N_{1} - N_{2}$ and $K = N_{1}N_{2} - K' = \tfrac14(S^{2} - D^{2}) - K'$, $V_{-1, 3/2} = (S - \tfrac32)^{2} + D^{2} - \tfrac14(S^{2} - D^{2}) + K' = \tfrac34(S - 2)^{2} + \tfrac54D^{2} - \tfrac34 + K'$, a sum of nonnegative terms and $-\tfrac34$, equal to $-\tfrac34$ exactly when $S = 2$, $D = 0$, $K' = 0$; by (ii) these minima have rank $3$, for example $\Phi = (1, 1)$, while the charge-breaking $(1, e_{02})$ with the same norms gives $\tfrac14$. $V_{1, 2} = (S - 2)^{2} + D^{2} + K \ge 0$, zero exactly when $N_{1} = N_{2} = 1$ and $K = 0$, hence $K' = 1$ and rank $4$; the aligned $(1, 1)$ gives $1$. A global minimum of a smooth function is stationary, so (i) applies to both. The rest of the BRST structure does not see $v$: $s^{2} = 0$ by (i), and the gauge-fixed action is BRST invariant, Fact \claimref{synthesis.imported-brst-kugo-ojima-strocchi-and-gribov}(i), its gauge-fixing term being $s\Psi$. In the 't Hooft-Feynman gauge ($\xi = 1$, canonical normalisation) the cross terms of $\tfrac12\langle D_{\mu}\Phi, D^{\mu}\Phi\rangle$ and $-\tfrac12\kappa_{ab}F^{a}F^{b}$ bilinear in $A$ and $h$ are $g\,\partial_{\mu}(A^{a\mu}\langle W_{a}, h\rangle)$ with $W_{a} = \rho(T_{a})v$, a total derivative, for every $v$; the squared masses of the vectors, the ghosts and the Goldstone modes are then the eigenvalues of $g^{2}\kappa^{-1}\Gamma$, $g^{2}\kappa^{-1}\Gamma$ and $g^{2}W\kappa^{-1}W^{\mathsf T}$ ($\Gamma_{ab} = \langle W_{a}, W_{b}\rangle$), the last because at a stationary point $\mathrm{Hess}\,V(v)$ annihilates the Goldstone directions (the identity of (i) with $\nabla V(v) = 0$), and $W\kappa^{-1}W^{\mathsf T}$ and $\kappa^{-1}W^{\mathsf T}W = \kappa^{-1}\Gamma$ have the same nonzero eigenvalues with multiplicity (the standard computation; the intended citation is named in the extraction note). For the free theory, diagonalise the symmetric mass matrix $M^{2}_{ab} = g_{a}g_{b}\langle\rho(T_{a})v, \rho(T_{b})v\rangle$ by an orthogonal change of the trace-orthonormal basis of the algebra: each generator with $M^{2} > 0$ has the Goldstone mode $\rho(T_{a})v$, these modes are mutually orthogonal, and each pair is the massive vector of Theorem \claimref{synthesis.quartet-homotopy-makes-gluon-and-massive-vector-physical-states-positive-at-every-number-of-quanta}(iii) with $\beta^{2} = Q(p)$, while the generators with $M^{2} = 0$ are gluon-like, its part (ii); the homotopy lemma, its part (i), then gives a positive physical space at every number of quanta, at either rank. The physical counts are $2(12 - r) + 3r$ vector modes and $8 - r$ scalars: $27 + 5$ at $r = 3$ (the stored count) and $28 + 4$ at $r = 4$. The check confirms the two decompositions on samples, the values at the minima and at the competing configurations, the ranks and the counts.

\emph{(iv)} In the setting of Theorem \claimref{synthesis.rung-field-on-the-lift-locus}, as used in Theorem \claimref{synthesis.rung-field-d-term-action-and-tree-masses-on-the-lift-locus}, the slot $U(2)$ acts by $\alpha \mapsto X\alpha$ and $\beta \mapsto \bar{X}\beta$, so $\bar{\beta} \mapsto X\bar{\beta}$ and both $\phi_{1} = \alpha$ and $\phi_{2} = \bar{\beta}$ carry the defining action; the rank of the Goldstone map and the invariants depend only on the real representation, which is the same for both. The check confirms on samples that under $(z_{1}, z_{2}) \mapsto z_{1} + e_{01}z_{2}$ the span$\{1, e_{12}\}$ part of $\widetilde{\phi}\psi$ is $\sum_{k}\bar{z}_{k}w_{k}$ and that left multiplication by $e_{12}$ is $\mathrm{diag}(e_{12}, -e_{12})$, the standard reading of a doublet as an even element. So $Q(c) = \alpha^{\mathsf T}\beta = \sum_{a}\overline{\bar{\beta}_{a}}\alpha_{a} = \phi_{2}^{\dagger}\phi_{1}$, $K = |Q(c)|^{2}$, and $\mathsf{N}^{2} - m^{2} = \tfrac14((N_{1} + N_{2})^{2} - (N_{1} - N_{2})^{2}) = N_{1}N_{2}$. On $L$, $K = H^{4} > 0$ and by (ii) $N_{1}N_{2} = H^{4} + K'$; so $K' = 0$, rank $3$, exactly on $\mathsf{N}^{2} - m^{2} = H^{4}$, which part (1) of the D-term theorem identifies with the charge-preserving orbits (the orbits with a circle stabiliser in Theorem \claimref{synthesis.rung-field-on-the-lift-locus}(1)), and rank $4$ inside; $\Phi \ne 0$ on $L$, so the rank is never $0$. Its part (2) puts the D-term vacuum on the boundary. The check confirms the rank on boundary and interior points of $L$ with $H^{2} = 1$.
